% paper/sections/03_levelset.tex
\section{Conservative Level Set (CLS) 法による界面追跡}
\label{sec:levelset}

\subsection{保存形移流方程式}

標準的な Level Set 法では質量保存性が保証されない問題があるため，
本稿では Conservative Level Set 法 \cite{OlssonKreiss2005} を採用する．
界面指標として，滑らか化ヘヴィサイド関数に相当する $\psi(\bm{x}, t)$ を用いる．
移流方程式は保存形で記述される：
\begin{equation}
  \pd{\psi}{t} + \bnabla\cdot(\psi\,\bu) = 0
  \label{eq:cls_advection}
\end{equation}
この式を解くことで，全領域積分 $\int \psi \, dV$（すなわち液相体積）の保存性が向上する．

\subsection{再初期化方程式（等方拡散形）}

移流によって $\psi$ のプロファイルが崩れるため，界面厚さを一定に保つ再初期化を行う．
CLS 法における再初期化は以下の圧縮・拡散方程式を仮想時間 $\tau$ 方向に解くことで実現される：
\begin{equation}
 \pd{\psi}{\tau} + \bnabla\cdot\bigl(\psi(1-\psi)\hat{\bm n}\bigr) = \bnabla\cdot\bigl(\varepsilon\,\bnabla\psi\bigr)
 \label{eq:cls_reinit_iso}
\end{equation}

ここで法線ベクトルは，幾何学的精度を保つため，符号付き距離関数 $\phi$ から $\hat{\bm n} = \bnabla \phi / |\bnabla \phi|$ として計算する．
右辺の拡散項係数 $\varepsilon$ は，数値的な界面厚さを決定する重要なパラメータである．
物理的な界面は厚さゼロの理想的な不連続面だが，これを離散格子上で直接表現しようとすると，ギブス現象による深刻な数値振動が発生する．
CLS法では，界面を有限の厚みを持つ滑らかな遷移層として扱うことで，この問題を回避する．

\begin{tcolorbox}[defbox, title={界面厚さの物理的・数値的トレードオフ}]
界面の厚み（拡散係数 $\varepsilon$）の決定は，数値的安定性と物理的忠実性のトレードオフを考慮する必要がある．

\begin{center}
    % Dummy figure for concept
    \includegraphics[width=0.9\textwidth]{placeholder_interface_profile.png}
    \textit{図X: 1次元における界面プロファイルの比較．(a) 理想的なステップ関数（物理的だが数値化困難），(b) 過度に厚い界面（数値的に安定だが不正確），(c) 適切な厚みの界面（両者の妥協点）．}
\end{center}

\begin{itemize}
    \item \textbf{厚すぎる界面（大きな $\varepsilon$）:}
    プロファイルが滑らかになり，数値計算は非常に安定する．しかし，物理的にシャープな界面からかけ離れてしまい，曲率計算の精度が低下したり，表面張力などの界面現象を不正確に捉えたりする原因となる．また，再初期化による質量保存誤差も増大しやすい．

    \item \textbf{薄すぎる界面（小さな $\varepsilon$）:}
    物理的な現実に近づくが，プロファイルが急峻になりすぎる．これは離散格子で解像しきれず，密度や粘性の急変化点で非物理的な振動を引き起こす．これは計算を不安定化させる主要な要因となる．
\end{itemize}

\textbf{妥協点：}
このトレードオフの標準的な解決策が，界面の厚さをグリッドサイズに比例させることである．
\begin{equation}
 \varepsilon = C_\varepsilon \, \Delta x_{\min} \quad (C_\varepsilon \approx 1.0 \sim 2.0)
 \label{eq:epsilon_scaling}
\end{equation}
ここで $\Delta x_{\min}$ は最小グリッド幅である．この設定により，界面は常に「数グリッド分」（典型的には $3 \sim 4$ グリッド）の厚さに保たれる．
これは，\textbf{グリッドが解像できる範囲で最もシャープな界面を維持する}，という合理的な妥協点であり，安定性と精度のバランスを取るための鍵となる．
\end{tcolorbox}

\begin{tcolorbox}[mybox, title={実装ガイド：再初期化ステップの運用}]
\textbf{反復回数と頻度：}
再初期化方程式 \eqref{eq:cls_reinit_iso} は，物理時間 $t$ とは独立した仮想時間 $\tau$ に関する偏微分方程式である．
実用上は，各物理時間ステップの終了時に，この方程式を数回（例えば $\Delta \tau$ の時間刻みで 3～5 ステップ）反復的に解くことで，崩れた $\psi$ のプロファイルを修復する．

\textbf{トレードオフ：}
反復回数が少なすぎると，$\psi$ がヘヴィサイド関数の形状を保てず，界面がぼやけたり，質量保存性が悪化したりする．
一方，反復回数が多すぎると，圧縮項（左辺第2項）の影響で界面位置がわずかに移動し，これもまた質量保存誤差の原因となる．
最適な反復回数や $\Delta \tau$ は問題設定に依存するため，予備計算を通じて調整することが望ましい．
一般的には，全領域で $\int \psi(1-\psi) dV$ が最小化される状態を目指す．
\end{tcolorbox}

\subsection{\texorpdfstring{$\psi$}{ψ} と \texorpdfstring{$\phi$}{φ} の関係}
\label{sec:levelset_mapping}

\begin{tcolorbox}[defbox]
\textbf{対応付け：}
Conservative Level Set 変数 $\psi$ と符号付き距離関数 $\phi$ は以下の関係にある：
\begin{equation}
  \psi(\bm{x},t) \equiv H_\varepsilon(\phi(\bm{x},t)), \qquad \bnabla\psi = \delta_\varepsilon(\phi)\,\bnabla\phi
\end{equation}
物性値補間 $\rho(\psi), \mu(\psi)$ には $\psi$ を直接用い，
曲率計算などの幾何学的操作には $\phi$ を用いる．

\vspace{1em}
\textbf{逆変換の必要性：}
CLS 法の核心は，移流計算は $\psi$ で行い，幾何量（特に再初期化に必要な法線 $\hat{\bm n}$）の計算は $\phi$ で行う点にある．
したがって，移流後の $\psi^{n+1}$ から対応する $\phi^{n+1}$ を求める「逆変換」が必要となる．
$\psi = H_\varepsilon(\phi)$ は非線形関数であるため，この逆変換は解析的に行えない．
各格子点において，$\psi_i$ の値に対して $H_\varepsilon(\phi_i) - \psi_i = 0$ を満たす $\phi_i$ を，ニュートン法などの反復解法を用いて数値的に求める必要がある．
$H_\varepsilon$ は単調増加関数であるため，解は一意に定まる．
\end{tcolorbox}
